\begin{tikzpicture}[line width=0.5pt]
% 坐标轴
\draw (-1.20,0) -- (6.30,0);
\node[anchor=west] at (6.28,0.00) {$K_X^{\perp}$};
\node[anchor=west] at (4.95,0.50) {$K_X<0$};
\node[anchor=west] at (4.95,-0.52) {$K_X>0$};
% 曲线锥 NE(X)
\fill[gray!20] (1.05,-0.05) -- (1.18,0.62) -- (1.78,1.32) -- (2.38,1.80)
  -- (2.72,2.04) -- (3.20,2.24) -- (3.72,2.20) -- (4.20,1.82) -- (4.50,1.12)
  -- (4.62,0.38) .. controls (4.58,-1.35) and (3.65,-2.55) .. (2.55,-2.58)
  .. controls (1.55,-2.60) and (1.00,-1.45) .. (1.05,-0.05) -- cycle;
\draw (1.05,-0.05) -- (1.18,0.62) -- (1.78,1.32) -- (2.38,1.80)
  -- (2.72,2.04) -- (3.20,2.24) -- (3.72,2.20) -- (4.20,1.82) -- (4.50,1.12)
  -- (4.62,0.38) .. controls (4.58,-1.35) and (3.65,-2.55) .. (2.55,-2.58)
  .. controls (1.55,-2.60) and (1.00,-1.45) .. (1.05,-0.05);
% 过原点的两条直线
\draw (0,0) -- (0.86,4.05);
\draw (0,0) -- (3.42,2.57);
\node[above] at (0.86,4.05) {$M=0$};
\node[above] at (3.36,2.57) {$M+rK_X=0$};
% 标注
\node[align=left,anchor=east] at (-0.05,2.95) {各种丰富\\除子 $M$};
\node[align=center] (ray) at (1.10,2.28) {此处有\\极射线};
\draw (1.50,2.18) -- (2.68,2.10);
\node at (2.75,-0.35) {$\overline{NE}(X)$};
\end{tikzpicture}
